\newpage
\phantomsection
\label{prf:ring-axioms-well-defined}

\begin{remark*}[Return]
\hyperref[def:ring]{Return to Definition: Ring}
\end{remark*}

\begin{remark*}[Note on proof type]
This file accompanies the Ring definition rather than a theorem. The
``proof'' here is a verification that the axiom set is internally
consistent (no axiom is derivable from the others in a way that
would create circular dependencies) and that the structure as given
is non-empty by example. It is not a proof of a theorem but a
justification that the definition is well-posed.
\end{remark*}

\begin{proposition*}[The Ring Axioms Are Well-Defined]
The axiom set \{(1),(2),(3),(5),(6),(7),(8),(9)\} of Halmos--Givant is
non-redundant and consistent: each axiom fails in some structure
satisfying all the others, and the integers $(\mathbb{Z},+,\cdot,-,0,1)$
witness consistency.
\end{proposition*}

\begin{proof}
TODO
\end{proof}
